\section{Related work}\label{sec:related}

\paragraph{Planted-record audits} The secret-sharer methodology plants random secrets and
measures their exposure \citep{carlini2019secret}. Privacy auditing with canaries estimates a
lower bound on the privacy parameter of a training algorithm from randomly included canaries
\citep{steinke2023privacy}, and copyright traps plant sequences in documents so that their
inclusion in training can be tested later \citep{meeus2024traps}. VOUCH
\citep{dang2026vouch} adds a paired control and betting e-processes
\citep{waudbysmith2024,ramdas2023game,shafer2021testing,grunwald2024safe} to certify the
outcome of one unlearning run with an exact null. None of these works bounds what a provider gains
by recognising the planted records. Theorem~\ref{thm:detect} covers an audit of this family
when its statistic is an average of bounded per-canary terms under randomised inclusion;
exposure ranks, abstaining guesses and AUC statistics require a separate argument.

\paragraph{Adversarial providers and verification} \citet{thudi2022necessity} argue that
approximate unlearning cannot be audited from weights alone. \citet{zhang2024verification}
show that a dishonest provider can pass existing families of verification strategies while
retaining the data; our bound quantifies this for planted-canary audits as a function of
canary detectability, and shows that it can be closed only at the price of power.
Cryptographic proofs of unlearning \citep{eisenhofer2025verifiable} attest that a committed
computation ran, which is complementary: they exclude deviation from the pipeline, while our
threat model is a pipeline that is run faithfully but treats records by content.
\citet{yu2025impossibility} show that retraining-based definitions are ill-posed for
multi-stage pipelines, one reason paired audits avoid a retrained reference.

\paragraph{Membership inference as an audit} Score-based attacks threshold the loss, its zlib
ratio \citep{carlini2021extracting} or Min-K\% Prob \citep{shi2024mink}; LiRA calibrates per
example with shadow models \citep{carlini2022lira}, and \citet{hayes2024inexact} show that
population-level attacks overstate unlearning. \citet{zhang2024mia} establish that
membership scores cannot prove training on given data without randomised inclusion.
\citet{duan2024membership} find attacks near chance when members and non-members share a
distribution, and \citet{das2024blind} show that classifiers that never query the model match
published attacks on shifted benchmarks. Our results add the unlearning-audit consequence:
an unmatched holdout accuses an honest provider, and a matched holdout is valid but weaker
than a paired test.

\paragraph{Evaluating unlearning} Gradient ascent, NPO \citep{zhang2024npo}, SimNPO
\citep{fan2025simnpo} and representation misdirection \citep{li2024wmdp} are evaluated on TOFU
\citep{maini2024tofu} and MUSE \citep{shi2024muse}. Their metrics have been shown to overstate
forgetting: unlearned knowledge returns after benign relearning \citep{hu2025jogging}, under
quantisation \citep{zhang2025quantization} and under lightweight attacks
\citep{lucki2024adversarial}, and forget-quality metrics track surface behaviour
\citep{wang2025towards,yuan2025closer,feng2025existing}; see
\citet{liu2025rethinking,xue2026survey} for surveys. \citet{li2026beliefs} show that
gradient-based unlearning shifts probability onto paraphrases of the target; our paraphrase
scores measure the audit side of the same phenomenon.
